\documentclass[12pt]{article}
\usepackage[utf8]{inputenc}
\usepackage[margin=1in]{geometry}
\usepackage{amsmath}
\usepackage{graphicx}
\usepackage{booktabs}
\usepackage{hyperref}
\usepackage[round]{natbib}

\title{Separating Local and Volume-Conducted Contributions to Thalamic Local Field Potentials Using Kernel Current Source Density Analysis in the Rat Somatosensory System}
\author{Author A$^{1}$, Author B$^{1,2}$, Author C$^{2}$, Author D$^{1}$ \\
\small $^{1}$Laboratory of Animal Physiology, Nencki Institute of Experimental Biology, Warsaw, Poland \\
\small $^{2}$Department of Neurophysiology, University of Warsaw, Warsaw, Poland}
\date{}

\begin{document}
\maketitle

\begin{abstract}
Local field potentials (LFPs) are widely used for population-level neural monitoring, but volume conduction causes distant high-amplitude sources to contaminate recordings at remote sites. In the rat whisker-barrel system, the large cortical evoked potential at approximately 10 ms post-stimulus may passively propagate into subcortical structures including the thalamus, confounding interpretation of thalamic signals. Here, we used kernel current source density (kCSD) analysis in 11 urethane-anesthetized Wistar rats with simultaneous recordings from barrel cortex and somatosensory thalamus (VPM/PoM) to quantify the cortical contribution to thalamic evoked potentials. Rolling correlation between thalamic and cortical EPs exceeded 0.5 at 10--12 ms post-stimulus, indicating substantial shared signal. kCSD source decomposition revealed that thalamic EPs reconstructed from local thalamic sources alone were significantly less negative at 10 ms than raw recorded EPs (permutation paired-test, $p < 0.05$), confirming cortical volume conduction contamination. The strong negative cortical field extended as a continuous stripe from cortex through the entire subcortical space to the thalamus. Critically, concurrent cortical recordings were not essential for reliable thalamic CSD estimation---thalamic-only recordings with sources assumed only in the thalamus yielded similar CSD profiles. Results were replicated across three probe technologies (Neuropixels 384-channel, NeuroNexus 8$\times$8, NeuroNexus 32-channel), confirming generalizability.
\end{abstract}

\section{Introduction}

Local field potentials (LFPs) reflect the summed synaptic and transmembrane currents of neural populations and are widely used to study circuit-level processing in the brain \citep{buzsaki2012origin, einevoll2013modelling}. However, because the extracellular medium is a volume conductor, electric fields generated by neuronal populations spread passively through tissue, potentially contaminating recordings at sites distant from the generating sources \citep{katzner2009local}. Estimates of the spatial extent of volume conduction range from hundreds of micrometers to several millimeters, depending on the geometry and amplitude of the source \citep{mitzdorf1985current}.

The rat whisker-barrel system is a well-characterized model for studying somatosensory processing across thalamocortical circuits \citep{petersen2007barrel, swadlow2003thalamocortical}. Whisker deflection produces a rapid sequence of neural responses: the ventral posteromedial thalamus (VPM) responds at approximately 3 ms post-stimulus, followed by barrel cortex activation at approximately 5 ms, with a large-amplitude negative evoked potential peaking in cortical layers 4/5 around 10 ms \citep{nicolelis1995somatotopic}. This cortical EP can exceed the thalamic signal by an order of magnitude, raising the question of whether a negative deflection observed in thalamic recordings at the same latency is genuinely local or partially volume-conducted from cortex.

Current source density (CSD) analysis can recover local current generators from LFP recordings by computing spatial derivatives of the potential field \citep{mitzdorf1985current}. The kernel CSD (kCSD) method extends classical CSD to handle arbitrary electrode geometries using a model-based kernel approach \citep{potworowski2012kcsd, leski2011inverse}. Importantly, kCSD allows selective reconstruction of field potentials from specific source subsets, enabling quantification of each structure's contribution to the recorded signal at any point.

Here, we applied kCSD to simultaneous cortical and thalamic recordings in the rat whisker system to: (1) quantify the cortical volume conduction contribution to thalamic EPs, (2) recover the truly local thalamic signal, and (3) determine whether concurrent cortical recordings are necessary for reliable thalamic CSD estimation.

\section{Methods}

\subsection{Animals and Surgical Preparation}

Eleven adult male Wistar rats (350--560 g) were anesthetized with urethane (1.5 mg/kg i.p., with 10\% supplemental doses as needed). Body temperature was maintained at 37--38$^\circ$C. All procedures complied with EU Directives 86/609/EEC and 2010/63/EU and were approved by the 1st Warsaw Local Ethics Committee.

\subsection{Recording Configurations}

Three probe configurations were used across the 11 animals (Table~\ref{tab:probes}). Cortical probes were tilted approximately 30$^\circ$ from vertical for perpendicular cortical entry; thalamic probes were inserted vertically. Electrode positions were verified histologically with cytochrome oxidase and Nissl staining.

\begin{table}[htbp]
\centering
\caption{Recording probe configurations.}
\label{tab:probes}
\begin{tabular}{llcc}
\toprule
Setup & Probe Type & Channels & Coverage \\
\midrule
1 & Neuropixels 1.0 & 384 & Continuous cortex-to-thalamus \\
2 & NeuroNexus 8$\times$8 ($\times$2) & 64 per probe & Separate cortex \& thalamus \\
3 & NeuroNexus 32-ch linear ($\times$2) & 32 per probe & Separate cortex \& thalamus \\
\bottomrule
\end{tabular}
\end{table}

\subsection{Stimulation}

Whisker stimulation was delivered via piezoelectric stimulator producing 0.1 mm horizontal deflections (1--2 ms duration, 20 V square wave). One hundred stimuli were delivered per session at 3--5 s pseudorandom intervals, controlled by Spike2 sequencer.

\subsection{kCSD Source Decomposition}

The kCSD method was applied to averaged evoked potential profiles \citep{potworowski2012kcsd}. Three analytical configurations were compared: (A) thalamic + cortical recordings with sources assumed only in the thalamus, (B) thalamic recordings alone with sources assumed only in the thalamus, and (C) full cortico-thalamic recordings with sources spanning the entire tissue block (Table~\ref{tab:configs}). After estimating the CSD, field potentials were selectively reconstructed from cortical sources only or thalamic sources only.

\begin{table}[htbp]
\centering
\caption{kCSD analytical configurations.}
\label{tab:configs}
\begin{tabular}{llll}
\toprule
Config & Recordings Used & Source Space & Purpose \\
\midrule
A & Thalamus + cortex & Thalamus only & Reference with cortical context \\
B & Thalamus only & Thalamus only & Test cortical recording necessity \\
C & Thalamus + cortex & Full block & Complete source reconstruction \\
\bottomrule
\end{tabular}
\end{table}

\subsection{Statistical Analysis}

Measured thalamic EPs at 10 ms post-stimulus were compared to reconstructed local-source-only EPs using permutation paired-tests. Rolling Pearson correlations (3-ms window) between thalamic and cortical EP waveforms were computed across all 11 rats.

\section{Results}

\subsection{Temporal Profile of Evoked Responses}

Whisker deflection produced a characteristic sequence of responses. The earliest thalamic response appeared at approximately 3 ms post-stimulus, consisting of single- and multi-unit activity in VPM/PoM. Cortical responses began at approximately 5 ms, with a large negative EP peak in layers 4/5 at approximately 10 ms. A negative deflection was also observed in the thalamic recording at the same 10-ms latency. Rolling correlation analysis between thalamic and cortical EPs across the 11 animals revealed correlation values exceeding 0.5 in the 10--12 ms window, indicating that a substantial portion of the signal was shared between the two recording sites at this critical time point.

\subsection{Cortical Volume Conduction Extends Through Subcortical Space}

Spatiotemporal EP profiles from Neuropixels recordings spanning cortex to thalamus revealed that the strong negative cortical field at 10 ms extended as a continuous stripe from the mid-cortical layers through the entire subcortical space to the thalamic recording sites. At the bottom of cortex, the field showed a brief polarity reversal (positive), but this was transient, and the negative field continued into thalamic territory. This spatial pattern is consistent with the cortical source generating a far-reaching dipole field that passively propagates to subcortical structures.

\subsection{kCSD Source Decomposition Reveals Cortical Contamination}

kCSD analysis decomposed the recorded signals into contributions from cortical and thalamic source populations. Reconstruction of field potentials from cortical sources alone demonstrated that cortical fields produced substantial signal at thalamic recording sites, particularly at 10 ms when the cortical EP peaked. Reconstruction from thalamic sources alone yielded a different waveform with smaller amplitude at 10 ms.

The critical quantitative test compared measured thalamic EPs with EPs reconstructed from local thalamic sources only. At 10 ms post-stimulus, the reconstructed values were systematically less negative than the measured values (permutation paired-test, $p < 0.05$). This confirms that the recorded thalamic negative deflection at 10 ms includes a significant component volume-conducted from the cortical source.

\subsection{Concurrent Cortical Recordings Are Not Essential}

A key practical finding was that configuration B (thalamic recordings only, with sources assumed only in the thalamus) yielded CSD profiles very similar to configuration A (which included cortical recordings). This demonstrates that concurrent cortical recordings are not necessary for reliable estimation of local thalamic CSD. The kCSD method, by constraining source estimation to the thalamic volume, effectively excludes cortical contributions even without explicitly measuring them.

\subsection{Generalization Across Probe Technologies}

Results were replicated across all three recording setups (Table~\ref{tab:probes}). The Neuropixels probe, with its 384 channels spanning cortex to thalamus continuously, provided the most complete spatial picture \citep{jun2017neuropixels}. The dual NeuroNexus 8$\times$8 configuration offered multi-shank coverage of both structures. The dual 32-channel linear configuration, while having the fewest channels, still yielded consistent kCSD source decomposition results. This confirms that the approach generalizes across probe technologies and is not dependent on specific electrode geometry.

\section{Discussion}

This study demonstrates that a substantial component of the negative evoked potential recorded in rat somatosensory thalamus around 10 ms post-whisker-stimulus is volume-conducted from the barrel cortex. The kCSD method successfully separates local and remote contributions, recovering the genuinely local thalamic signal and quantifying the cortical contamination.

\subsection{Implications for LFP Interpretation}

Our findings have direct implications for the interpretation of thalamic LFP recordings in the whisker-barrel system and beyond \citep{buzsaki2012origin, einevoll2013modelling}. The cortical barrel field generates one of the largest evoked potentials in the rodent brain, and our results show that this field extends through the entire subcortical space to the level of the thalamus. Studies interpreting thalamic LFP waveforms at latencies overlapping with the cortical EP (approximately 8--15 ms) should account for this contamination, either by applying CSD analysis or by interpreting only components that are clearly separable in time from the cortical peak \citep{steriade2001thalamocortical}.

The earliest thalamic response at approximately 3 ms precedes cortical activation and is therefore free from cortical contamination. This provides a clean window for assessing purely thalamic processing. The volume conduction concern is greatest in the 8--15 ms range where cortical and thalamic responses overlap temporally.

\subsection{Practical Implications for Experimental Design}

The finding that concurrent cortical recordings are not essential for reliable thalamic CSD estimation is practically important. It means that researchers studying thalamic processing need not always instrument the cortex simultaneously, provided they apply appropriate CSD analysis to their thalamic recordings \citep{potworowski2012kcsd, leski2011inverse}. This simplifies experimental design and reduces surgical complexity, particularly in chronic recording preparations.

\subsection{Broader Relevance}

Volume conduction is a general concern for all LFP recordings, not limited to the somatosensory system. Hippocampal oscillations, for example, can propagate widely, and human intracranial recordings face similar contamination issues \citep{katzner2009local, mitzdorf1985current}. The kCSD-based source decomposition approach demonstrated here could be applied to any multi-structure recording configuration where volume conduction is a concern.

\subsection{Limitations}

The study was conducted under urethane anesthesia, which may alter the amplitude and temporal profile of both cortical and thalamic evoked responses. The permutation test at 10 ms represents a single time-point comparison; future work should extend this to full waveform statistical analysis. The kCSD method assumes homogeneous tissue conductivity, which is an approximation.

\section{Conclusion}

Using kernel CSD analysis across three probe technologies in 11 rats, we demonstrate that thalamic evoked potentials in the rat whisker system contain a significant volume-conducted contribution from barrel cortex, particularly at 10 ms post-stimulus when the cortical EP peaks. The kCSD method successfully recovers the truly local thalamic signal and does not require concurrent cortical recordings. These findings underscore the importance of CSD analysis for accurate interpretation of subcortical LFP recordings and provide a validated methodological framework applicable to any multi-structure recording configuration.

\bibliographystyle{plainnat}
\bibliography{references}

\end{document}
